% Part: first-order-logic
% Chapter: completeness
% Section: introduction

\documentclass[../../../include/open-logic-section]{subfiles}

\begin{document}

\iftag{FOL}
      {\olfileid[hi]{fol}{com}{int}}
      {\olfileid[hi]{pl}{com}{int}}

\olsection{परिचय}

पूर्णता प्रमेय तर्कशास्त्र के सबसे मूलभूत परिणामों में से एक है। इसके दो
सूत्रीकरण हैं, जिनकी समतुल्यता हम सिद्ध करेंगे। पहले सूत्रीकरण में यह
अर्थवैज्ञानिक निष्पत्ति और हमारी !!{derivation}-प्रणाली के संबंध के बारे में
एक मूलभूत बात कहता है: यदि कोई !!a{sentence}~$!A$, कुछ !!{sentence}s
$\Gamma$ से तार्किक रूप से निकलता है, तो ऐसी !!a{derivation} भी होती है जो
$\Gamma
\Proves !A$ स्थापित करती है। अतः !!{derivation}-प्रणाली उन बातों को
सिद्ध किए बिना, जो वास्तव में निष्पन्न नहीं होतीं, जितनी प्रबल हो सकती है
उतनी प्रबल है।

दूसरे सूत्रीकरण में इसे प्रतिरूप-अस्तित्व परिणाम के रूप में कहा जा सकता है:
!!{sentence}s का प्रत्येक संगत समुच्चय संतुष्टनीय है। संगतता एक
प्रमाण-सैद्धांतिक धारणा है: इसका अर्थ है कि हमारी !!{derivation}-प्रणाली कुछ
विशेष प्रकार की !!{derivation}s उत्पन्न नहीं कर सकती। किंतु केवल इस कारण कि
कुछ विशेष प्रकार की !!{derivation}s~$\Gamma$ से नहीं मिलतीं, कौन कह सकता है
कि ऐसी
\iftag{FOL}{!!a{structure}~$\Struct{M}$}{!!{valuation}{~$\pAssign{v}$}
अवश्य है जिसमें $\iftag{FOL}{\Sat{M}}{\pSat{v}}{\Gamma}$}? पूर्णता प्रमेय
पहली बार सिद्ध होने से पहले---वस्तुतः आज की हमारी !!{derivation}-प्रणालियाँ
बनने से भी पहले---महान जर्मन गणितज्ञ डेविड हिल्बर्ट का मत था कि गणितीय
सिद्धांतों की संगतता उन वस्तुओं के अस्तित्व की गारंटी देती है जिनके बारे में
वे सिद्धांत हैं। गॉटलोब फ़्रेगे को लिखे एक पत्र में उन्होंने इसे इस प्रकार
कहा:
\begin{quote}
  यदि मनमाने ढंग से दिए गए अभिगृहीत अपने सभी परिणामों सहित एक-दूसरे का
  खंडन नहीं करते, तो वे सत्य हैं और अभिगृहीतों द्वारा परिभाषित वस्तुएँ
  विद्यमान हैं। मेरे लिए सत्य और अस्तित्व का निकष यही है।
\end{quote}
फ़्रेगे ने इसका तीव्र विरोध किया। पूर्णता प्रमेय का दूसरा सूत्रीकरण दिखाता
है कि कम-से-कम इस अर्थ में हिल्बर्ट सही थे कि यदि अभिगृहीत संगत हों, तो ऐसी
\emph{कोई}
\iftag{FOL}{!!{structure}}{!!{valuation}} विद्यमान होती है जो उन सबको सत्य
बनाती है।

पूर्णता प्रमेय---या अधिक ठीक कहें तो उसका प्रमाण---के महत्त्व के ये ही कारण
नहीं हैं। उसके कई महत्त्वपूर्ण परिणाम हैं, जिनमें से कुछ पर हम अलग से चर्चा
करेंगे। उदाहरणार्थ, $\Gamma \Proves !A$ दिखाने वाली प्रत्येक
!!{derivation} परिमित होती है और इसलिए~$\Gamma$ के केवल परिमिततः अनेक
!!{sentence}s का उपयोग कर सकती है। अतः पूर्णता प्रमेय से मिलता है कि यदि
$!A$,~$\Gamma$ का परिणाम है, तो वह~$\Gamma$ के किसी परिमित उपसमुच्चय का
परिणाम पहले ही है। इसे \emph{संहतता} कहते हैं। समतुल्यतः, यदि $\Gamma$ का
प्रत्येक परिमित उपसमुच्चय संगत है, तो स्वयं $\Gamma$ भी संगत होना चाहिए।

यद्यपि संहतता प्रमेय, !!{derivation}s के रास्ते होकर, पूर्णता प्रमेय से
निकलता है, फिर भी पूर्णता प्रमेय के \emph{प्रमाण} से उसे सीधे स्थापित करना
भी संभव है। वह प्रमाण एक विशेष गुण---संगतता---वाला !!{sentence}s का समुच्चय
लेता है और उसी से कुछ विशेष गुणों वाली !!a{structure} बनाता है (इस प्रसंग
में वह संरचना उस समुच्चय को संतुष्ट करती है)। लगभग उसी रचना से संहतता को सीधे
स्थापित किया जा सकता है: संगत समुच्चयों के स्थान पर !!{sentence}s के
``परिमिततः संतुष्टनीय'' समुच्चयों से आरंभ करें।
\iftag{FOL}{यह रचना अन्य परिणाम भी देती है; उदाहरणार्थ, !!{sentence}s के
  प्रत्येक संतुष्टनीय समुच्चय का कोई परिमित अथवा !!{denumerable}s प्रतिरूप
  होता है। (इस परिणाम को लोवेनहाइम--स्कोलेम (L\"owenheim--Skolem) प्रमेय
  कहते हैं।) सामान्यतः !!{structure}s को !!{sentence}s के समुच्चयों से रचने
  की विधि तर्कशास्त्र में और कभी-कभी दर्शन में भी प्रचलित है।}{}

\end{document}
